Agradeço, em primeiro lugar, a Deus por me permitir chegar até aqui, e a todas as pessoas que, por meio de suas orações e de sua presença, contribuíram de alguma forma para a minha caminhada.

À minha avó, Maria de Lourdes, que, para mim, foi muito mais do que uma avó. Sempre fez o possível e o impossível para que eu continuasse meus estudos. Infelizmente, não pôde ver a conclusão desta etapa, mas tenho certeza de que, onde quer que esteja, está orgulhosa.

Aos meus pais, Rodrigo e Reylla, pelo apoio incondicional e por me guiarem pelos caminhos certos da vida. Essa conquista não é só minha, mas de vocês também.

À minha prima Flavia, que é como uma irmã para mim. Entre brigas e risadas, sempre esteve ao meu lado nos momentos mais difíceis. E à sua mãe, Fátima, que, junto dela, alegrava até os dias mais escuros, mas que, assim como minha avó, deixa saudades que não cabem em palavras.

Aos amigos que a graduação me deu, companheiros nos mais diversos desafios, desde as chuvas até as matérias que pareciam impossíveis. Um agradecimento especial ao Helcio, que esteve ao meu lado na construção deste trabalho, e aos demais: Pedro ``Albergue'', ``Magé'', Mariana, Pedro Guedes, Hélio e Davi.

Por fim, agradeço ao Clube de Regatas do Flamengo, que me proporcionou incontáveis alegrias ao longo desses anos e que, à sua maneira, também fez parte desta jornada.
